\section{Síntesis y ampliación}

\subsection{Un hilo conductor que atraviesa toda la clase: diseñar primero la interfaz de memoria}

Esta clase puede resumirse en cuatro pasos de razonamiento. Primero, cada modelo autoregresivo tiene un estado interpasos. Segundo, la forma del estado determina la interfaz de acceso: una base de datos de tokens permite recuperación exacta, mientras que un autómata finito soporta compresión continua. Tercero, la resolución de las unidades de datos determina qué interfaz es natural; la atención requiere tokens efectivos y los SSM manejan mejor flujos de alta resolución que requieren compresión previa. Cuarto, H-Net demuestra una combinación: capas bajas aprenden chunks y las capas altas realizan modelado de relaciones más fuertes sobre los chunks.

\begin{importantbox}{Cuatro frases para recordar}
\begin{enumerate}
  \item El núcleo de Mamba es la combinación sistemática de expansión de estado, selectividad y escaneo eficiente.
  \item La verdadera ventaja de la atención es el acceso aleatorio con resolución de tokens, no solo una fórmula $QK^\top$.
  \item La tokenización/patchificación cambia el problema de modelado, no solo acorta las secuencias.
  \item Los sistemas más prometedores suelen ser híbridos conscientes de la resolución, no un primitivo único que domine todas las capas.
\end{enumerate}
\end{importantbox}

\subsection{Tres límites de evidencia}

Primero, perplexidad, BPB y pérdida de preentrenamiento son métricas locales que no pueden reemplazar directamente recuperación, razonamiento, control y ciencia downstream. Segundo, los resultados con bytes/DNA muestran la interacción entre arquitectura y resolución de tokens, pero no se puede concluir que todas las modalidades originales sean ganadas por SSM. Tercero, los límites aprendidos de H-Net y el ablation de etapas externas apoyan la hipótesis de compresión, pero no identifican de manera única que el "estado finito" sea la causa causal del beneficio.

\subsection{Lista de verificación para diseño de sistemas}

\begin{enumerate}
  \item \textbf{Definir estado:} ¿Qué se guarda después de cada paso generativo? ¿El estado crece con el contexto?
  \item \textbf{Definir recuperación:} ¿La tarea requiere mirar atrás exactamente en alguna unidad histórica o basta con un resumen suficiente?
  \item \textbf{Revisar granularidad:} ¿Tienen los tokens/patches/eventos actuales significado independiente? ¿Las fronteras deberían ser dependientes del contexto?
  \item \textbf{Ajustar a la capacidad de cómputo:} ¿La complejidad teórica, el tráfico HBM, la utilización del kernel y el comportamiento por lote son consistentes?
  \item \textbf{Realizar ablaciones:} Comparar tipos de estado y proporción de capas solo después de ajustar parámetros, FLOPs, datos y receta de entrenamiento.
  \item \textbf{Mantener límites:} No exagerar una curva de preentrenamiento como conclusión universal sobre inteligencia.
\end{enumerate}

\subsection{Preguntas abiertas}

Futuras preguntas valiosas para seguir incluyen: ¿puede el chunking dinámico escalar hasta la frontera; ¿puede la memoria jerárquica conservar simultáneamente un resumen grueso y detalles recuperables; ¿cómo hacer trazado causal de la dependencia difusa de los SSM; ¿puede el algoritmo de aprendizaje liberarse de las limitaciones de memoria de backprop para dependencias ultra largas; ¿proveerá el hardware óptimos distintos para la verdadera recurrencia; y en los híbridos, si la capa de atención debe activarse por ciclo fijo, enrutamiento dinámico o condiciones de tarea.

\subsection{Lecturas adicionales}

\begin{itemize}
  \item Albert Gu, \href{https://goombalab.github.io/blog/2025/tradeoffs/}{On the Tradeoffs of SSMs and Transformers}: origen conceptual de larga duración de esta clase.
  \item Gu y Dao, \href{https://arxiv.org/abs/2312.00752}{Mamba}; Dao y Gu, \href{https://arxiv.org/abs/2405.21060}{Mamba-2 / Duality de espacios de estado estructurados}.
  \item Lahoti et al., \href{https://arxiv.org/abs/2603.15569v1}{Mamba-3 v1}: actualización de SSM orientada a inferencia publicada un mes antes de la clase.
  \item Hwang, Wang y Gu, \href{https://arxiv.org/abs/2507.07955v2}{H-Net v2}: método completo de chunking dinámico y resultados de escalado.
  \item Wang et al., \href{https://arxiv.org/abs/2401.13660}{MambaByte}: SSM selectivo a nivel de byte; Shah et al., \href{https://arxiv.org/abs/2602.10603v3}{dnaHNet v3}: jerarquía genómica en el snapshot de la clase.
\end{itemize}

\end{document}
